% BTVN_B5.tex — bien tap tu ba de TDMXX09, TDMXX10, TDMXX11.
% Do thi duoc ve lai bang TikZ. Bien tap theo phong cach BTVN_01_B1.tex.
\documentclass[12pt,a4paper]{article}
\usepackage[left=15mm,right=15mm,top=14mm,bottom=18mm,footskip=9mm]{geometry}
\usepackage{iftex}
\ifPDFTeX
  \usepackage[T5]{fontenc}\usepackage[utf8]{inputenc}\usepackage{mathptmx}
\else
  \usepackage{fontspec}
  \IfFontExistsTF{Times New Roman}{\setmainfont{Times New Roman}}{%
    \IfFontExistsTF{Liberation Serif}{\setmainfont{Liberation Serif}}{\setmainfont{Latin Modern Roman}}}
\fi
\usepackage{amsmath,amssymb,fix-cm}
\usepackage{tikz}
\usetikzlibrary{arrows.meta}
\usepackage{tcolorbox,tabularx,fancyhdr,lastpage}
\definecolor{sectionpurple}{RGB}{155,42,143}
\definecolor{notebackground}{RGB}{229,240,245}
\definecolor{noteborder}{RGB}{255,55,55}
\definecolor{grblue}{RGB}{37,77,194}
\setlength{\parindent}{0pt}
\setlength{\parskip}{0pt}
\setlength{\headheight}{15pt}
\pagestyle{fancy}\fancyhf{}
\renewcommand{\headrulewidth}{0pt}\renewcommand{\footrulewidth}{0.4pt}
\fancyfoot[L]{\fontsize{8.5}{10}\selectfont Biên soạn: \textbf{Nguyễn Anh Dũng}}
\fancyfoot[R]{\fontsize{8.5}{10}\selectfont Trang \thepage/\pageref{LastPage} -- BTVN B5}
\newtcolorbox{notebox}{colback=notebackground,colframe=noteborder,
 boxrule=.65pt,arc=1.5pt,left=5pt,right=5pt,top=2pt,bottom=2pt,
 before skip=2pt,after skip=2pt}
\newcommand{\question}[2][TDM31]{%
  \par\addvspace{2pt}\noindent\textbf{\textcolor{sectionpurple}{Câu #2:} \textcolor{noteborder}{[#1]}}\enspace}
\newcommand{\answerlabel}[1]{\textbf{\textcolor{sectionpurple}{#1.}}}
\newcommand{\fourchoices}[4]{%
 \par\nobreak\vspace{1pt}\begin{tabular*}{\linewidth}{@{}l@{\extracolsep{\fill}}lll@{}}
 \answerlabel{A} #1 & \answerlabel{B} #2 & \answerlabel{C} #3 & \answerlabel{D} #4
 \end{tabular*}\par}
\newcommand{\twochoices}[4]{%
 \par\nobreak\vspace{1pt}\begin{tabularx}{\linewidth}{@{}>{\raggedright\arraybackslash}X>{\raggedright\arraybackslash}X@{}}
 \answerlabel{A} #1 & \answerlabel{B} #2 \\[1pt]
 \answerlabel{C} #3 & \answerlabel{D} #4
 \end{tabularx}\par}
\newcommand{\tfbox}[4]{%
\begin{notebox}\renewcommand{\arraystretch}{1.17}\setlength{\tabcolsep}{4pt}
\begin{tabularx}{\linewidth}{|>{\bfseries\color{sectionpurple}}c|>{\raggedright\arraybackslash}X|c|c|}
\hline\multicolumn{2}{|c|}{\textbf{\textcolor{sectionpurple}{Mệnh đề}}}&
\textbf{\textcolor{sectionpurple}{Đ}}&\textbf{\textcolor{sectionpurple}{S}}\\\hline
 a & #1 & $\bigcirc$ & $\bigcirc$ \\ \hline
 b & #2 & $\bigcirc$ & $\bigcirc$ \\ \hline
 c & #3 & $\bigcirc$ & $\bigcirc$ \\ \hline
 d & #4 & $\bigcirc$ & $\bigcirc$ \\ \hline
\end{tabularx}\end{notebox}}
\newcommand{\shortanswer}{%
 \par\nobreak\hfill\begin{tikzpicture}[baseline=-.7ex]
 \node[draw=noteborder,fill=notebackground,line width=.6pt,rounded corners=1pt,
 minimum width=23mm,minimum height=5mm,inner sep=0pt] {};
 \end{tikzpicture}\par}
\newcommand{\examtitle}[2]{%
 \par\begin{center}\textbf{ĐỀ #1}\\[-1pt]\textbf{#2}\\[-1pt]
 {\fontsize{10.5}{12}\selectfont\itshape Form đề chuẩn bám sát BGD2026 -- Thời gian làm bài \textbf{90 phút}}\end{center}\vspace{-7pt}}
\newcommand{\partone}[1]{\par\addvspace{3pt}\textbf{Phần I. Câu trắc nghiệm nhiều phương án lựa chọn.} Thí sinh trả lời từ câu 1 đến câu #1. Mỗi câu chỉ chọn một phương án.\par}
\newcommand{\parttwo}[2]{\par\addvspace{4pt}\textbf{Phần II. Câu trắc nghiệm đúng sai.} Thí sinh trả lời từ câu #1 đến câu #2, trong mỗi ý a), b), c), d) chọn đúng hoặc sai.\par}
\newcommand{\partthree}[2]{\par\addvspace{4pt}\textbf{Phần III. Câu trắc nghiệm trả lời ngắn.} Thí sinh trả lời từ câu #1 đến câu #2.\par}
\newcommand{\finishExam}{\par\vspace{3pt}{\centering\textbf{\textcolor{sectionpurple}{--- Hết ---}}\par}}
% Hinh ve duoc chuan hoa theo do thi vector goc PDF (TDMXX10) va doi chieu hinh raster goc (TDMXX09, TDMXX11).
% Hinh TikZ doc lap, khong can tai tep hinh ben ngoai.
\newcommand{\figbox}[1]{%
\begin{notebox}\centering
\begin{tikzpicture}[x=.43cm,y=.43cm,>=Stealth,every node/.style={font=\fontsize{10}{11}\selectfont}]
 #1
\end{tikzpicture}\end{notebox}}
\newcommand{\axes}[4]{%
 \draw[thin,-{Stealth[length=4pt]}] (#1,0)--(#2,0) node[below] {$x$};
 \draw[thin,-{Stealth[length=4pt]}] (0,#3)--(0,#4) node[right] {$y$};
 \node[below left] at (0,0) {$O$};}
\newcommand{\dotat}[2]{\fill (#1,#2) circle (1.3pt);}
\newcommand{\vguide}[2]{\draw[thin,densely dashed] (#1,0)--(#1,#2);\dotat{#1}{0}\dotat{#1}{#2}}
\begin{document}
\fontsize{11.5}{14}\selectfont
\examtitle{TDMXX09}{CƠ BẢN VỀ TIỆM CẬN XIÊN CỦA ĐỒ THỊ HÀM SỐ}
\partone{12}
\question[TDM21]{1} Tiệm cận xiên của đồ thị hàm số $y=x+2+\dfrac{1}{x-1}$ là
\fourchoices{$y=x+2$}{$x=-2$}{$x=1$}{$y=2$}
\question[TDM21]{2} Tiệm cận xiên của đồ thị hàm số $y=\dfrac{x^2-3x}{x}$ là
\fourchoices{$y=x+3$}{$y=2$}{$y=x-3$}{$x=0$}
\question[TDM21]{3} Tiệm cận xiên của đồ thị hàm số $y=\dfrac{x^3+1}{x^2-2}$ là
\fourchoices{$y=x^2-2$}{$y=\sqrt2$}{$x=\sqrt2$}{$y=x$}
\question[TDM21]{4} Tiệm cận xiên của đồ thị hàm số $y=x+3+\dfrac{x-2}{x-1}$ là
\fourchoices{$y=x+3$}{$y=x+4$}{$x=4$}{$y=2x+3$}
\question[TDM21]{5} Tiệm cận xiên của đồ thị hàm số $y=x-2+\dfrac{x^2}{x+1}$ là
\fourchoices{$y=x-2$}{$y=x-3$}{$y=2x-3$}{$y=x+6$}
\question[TDM31]{6} Số đường tiệm cận của đồ thị hàm số $y=\dfrac{x^3+10}{x^2-3x+2}$ là
\fourchoices{$2$}{$1$}{$0$}{$3$}
\question[TDM31]{7} Cho hàm số $y=x+1+\dfrac{1}{\sqrt{2-x}}+e^x$ có đồ thị $(C)$. Đường tiệm cận xiên của $(C)$ là
\fourchoices{$y=x$}{$y=2x+1$}{$y=x+1$}{$y=-x-1$}
\question[TDM31]{8} Cho hàm số $y=x-\sqrt{x^2-4x}$ có đồ thị $(C)$. Đường tiệm cận xiên của $(C)$ là
\fourchoices{$y=2x+1$}{$y=2x-2$}{$y=x+2$}{$y=-2x+4$}
\question[TDM31]{9} Biết rằng tiệm cận xiên và tiệm cận đứng của đồ thị hàm số $y=\dfrac{x^2-4x+2}{x-1}$ có phương trình lần lượt là $y=x-a$ và $x=b$. Giá trị $T=ab$ bằng bao nhiêu?
\fourchoices{$6$}{$3$}{$-12$}{$4$}
\question[TDM31]{10} Tổng số các đường tiệm cận của đồ thị hàm số $y=\dfrac{x^2}{x-1}-\sqrt{x^2+2x}$ là
\fourchoices{$2$}{$3$}{$1$}{$0$}
\question[TDM31]{11} Cho hàm số $y=f(x)=\dfrac{ax^2+(2a-1)x+b}{x-a}$ có đồ thị $(C)$. Biết $(C)$ có tiệm cận đứng là $x=2$. Đường tiệm cận xiên của $(C)$ có phương trình là
\fourchoices{$y=2x+7$}{$y=x-2$}{$y=2x-1$}{$y=2x+3$}
\question[TDM31]{12} Cho hàm số $y=f(x)$ có đồ thị $(C)$ và thỏa mãn $\displaystyle\lim_{x\to\pm\infty}(2f(x)-4x)=2$. Tiệm cận xiên của đồ thị $(C)$ là
\fourchoices{$y=4x+2$}{$y=2x+1$}{$y=2x-2$}{$y=2x+2$}
\parttwo{13}{16}
\question[TDM31]{13} Cho hàm số $y=f(x)$ có đồ thị $(C)$. Hỏi mệnh đề nào dưới đây là đúng hay sai?
\tfbox{Đường tiệm cận xiên của $(C)$ có thể cắt $(C)$.}
{Đường tiệm cận xiên của $(C)$ không thể có điểm chung với $(C)$.}
{Đồ thị $(C)$ có tối đa ba đường tiệm cận xiên.}
{Đồ thị $(C)$ luôn có ít nhất một đường tiệm cận xiên.}
\question[TDM31]{14} Cho hàm số $y=f(x)$ có đồ thị $(C)$ là hai nhánh như hình vẽ. Hỏi mệnh đề nào dưới đây là đúng hay sai?
\figbox{\axes{-5.6}{8.73}{-6.04}{8.96}
\draw[noteborder,thick] (-3.76,8.64) .. controls (-2.47,6.84) and (-1.24,4.53) .. (-0.2,4.51) .. controls (0.49,4.51) and (0.76,6.76) .. (0.84,8.98);
\draw[noteborder,thick] (1.07,-5.82) .. controls (1.13,-3.67) and (1.27,-1.71) .. (2.07,-1.58) .. controls (3.07,-1.2) and (4.47,-4.71) .. (5.44,-5.64);
\draw[black,thin] (-3.98,8.62)--(5.62,-5.69);
\draw[black,thin] (0.93,8.98)--(0.93,-6.04);
\draw[black,thin,densely dashed] (0,1.4)--(0.93,1.4);
\dotat{0}{0} \dotat{0}{1.4} \dotat{0}{2.82} \dotat{0.93}{0} \dotat{0.93}{1.4} \node[left] at (0,1.4) {$3$};
\node[left] at (0,2.82) {$6$};
\node[below] at (0.93,0) {$2$};
\node[right] at (-2.47,7.56) {$(C)$};
}
\tfbox{Đồ thị $(C)$ có đường tiệm cận đứng là $x=2$.}
{Đồ thị $(C)$ có đường tiệm ngang là $y=3$.}
{Đồ thị $(C)$ có đường tiệm cận xiên là $y=-\dfrac32x+6$.}
{Nếu hàm số có dạng $y=f(x)=\dfrac{ax^2+bx+c}{x+n}$ thì $abn=-27$.}
\question[TDM32]{15} Cho hàm số $y=f(x)$ có đồ thị $(C)$. Biết hàm số thỏa mãn các giới hạn $\displaystyle\lim_{x\to+\infty}(f(x)-3x)=2$ và $\displaystyle\lim_{x\to-\infty}(f(x)-2x)=3$. Hỏi mệnh đề nào dưới đây là đúng hay sai?
\tfbox{Đồ thị $(C)$ không có tiệm cận xiên.}
{Đồ thị $(C)$ có đường tiệm cận xiên là $y=3x+2$.}
{Đồ thị $(C)$ có đường tiệm cận xiên là $y=2x+3$.}
{Giao của hai đường tiệm cận xiên của $(C)$ có tọa độ là $(1;5)$.}
\question[TDM32]{16} Cho hàm số $y=f(x)=\dfrac{ax+2}{x+d}$ có đồ thị $(C)$. Biết tiệm cận đứng và tiệm cận ngang của $(C)$ có phương trình lần lượt là $x=2$ và $y=3$. Hỏi mệnh đề nào dưới đây là đúng hay sai?
\tfbox{Giá trị của $a$ bằng $3$.}
{Giá trị của $d$ bằng $2$.}
{Đồ thị hàm số $y=xf(x)$ có đường tiệm cận xiên là $y=3x+8$.}
{Để đồ thị hàm số $y=\dfrac{x+b}{cx+1}f(x)$ có đường tiệm cận xiên là $y=3x+5$ thì $b+c=-1$.}
\partthree{17}{22}
\question[TDM31]{17} Cho hàm số $y=f(x)=\dfrac{x^3+4x}{x^2-1}$ có đồ thị $(C)$. Hãy tính diện tích hình phẳng giới hạn bởi ba đường tiệm cận của $(C)$ với trục hoành $Ox$ (làm tròn kết quả đến hàng đơn vị)?
\shortanswer
\question[TDM31]{18} Biết rằng tiệm cận xiên và tiệm cận đứng của đồ thị hàm số $y=\dfrac{x^2+ax+10}{x+b}$ có phương trình lần lượt là $y=x-2$ và $x=3$. Giá trị biểu thức $T=a+b$ bằng bao nhiêu?
\shortanswer
\question[TDM31]{19} Cho hàm số $y=2x-\sqrt{x^2-6x+1}$ có đồ thị $(C)$. Biết $(C)$ có hai đường tiệm cận xiên là $y=ax+b$ và $y=cx+d$. Giá trị $T=ad+bc$ bằng bao nhiêu?
\shortanswer
\question[TDM31]{20} Cho hàm số $y=f(x)$ có đồ thị $(C)$. Biết với hai số thực $a,b$ ta có
$\displaystyle\lim_{x\to+\infty}(f(x)-ax+2b)=0$ và $\displaystyle\lim_{x\to-\infty}(f(x)-(3b-2)x+a+1)=0$.
Để có đúng một đường tiệm cận xiên thì $T=3a+b$ bằng bao nhiêu?
\shortanswer
\question[TDM32]{21} Biết tiệm cận xiên của đồ thị hàm số $y=f(x)=\dfrac{x^2+2mx+5}{x-1}$ cách đường thẳng $y=x$ một khoảng bằng $\sqrt2$. Hãy tính tổng tất cả các giá trị thực của $m$ thỏa mãn bài toán?
\shortanswer
\question[TDM32]{22} Biết tiệm cận xiên của đồ thị hàm số $y=f(x)=\dfrac{(a+b)x^3+4x^2+ax+3}{(2a+b)x-1}$ có phương trình $y=2x+c$. Giá trị biểu thức $T=2a+b+c$ bằng bao nhiêu?
\shortanswer
\finishExam
\clearpage
\examtitle{TDMXX10}{KHẢO SÁT HÀM SỐ BẬC BA}
\partone{12}
\question[TDM31]{1} Cho hàm số $f(x)=ax^3+bx^2+cx+d$ có đồ thị như hình vẽ bên dưới. Nhận xét nào dưới đây \textbf{đúng}?
\figbox{\axes{-9.47}{10.63}{-5.74}{8.53}
\draw[grblue,thick] (-6.84,-3.53) .. controls (-5.89,2.11) and (-5.11,4.95) .. (-3.89,4.95) .. controls (-2.11,4.95) and (-0.11,-3.37) .. (1.89,-3.37) .. controls (3.68,-3.37) and (4.74,3.32) .. (5.32,5.95);
\draw[black,thin,densely dashed] (-3.89,0)--(-3.89,4.95);
\draw[black,thin,densely dashed] (1.89,0)--(1.89,-3.37);
\dotat{-3.89}{4.95} \dotat{-3.89}{0} \dotat{1.89}{-3.37} \dotat{1.89}{0} \node[right] at (5.63,6.16) {$f(x)$};
}
\twochoices{$a>0;\ b>0;\ c<0;\ d<0$.}{$a>0;\ b<0;\ c<0;\ d<0$.}
{$a>0;\ b>0;\ c>0;\ d>0$.}{$a<0;\ b<0;\ c>0;\ d<0$.}
\question[TDM31]{2} Cho hàm số $f(x)=ax^3+bx^2+cx+d$ có đồ thị như hình vẽ bên dưới. Nhận xét nào dưới đây \textbf{đúng}?
\figbox{\axes{-9.47}{10.58}{-5.63}{7.26}
\draw[grblue,thick] (-5.79,-4.32) .. controls (-5.11,-0.05) and (-4.16,4.21) .. (-2.89,4.21) .. controls (-1.21,4.21) and (0.79,-4.11) .. (2.74,-4.11) .. controls (4.53,-4.11) and (5.89,3.11) .. (6.42,6.79);
\draw[black,thin,densely dashed] (-2.89,0)--(-2.89,4.21);
\draw[black,thin,densely dashed] (2.74,0)--(2.74,-4.11);
\dotat{-2.89}{4.21} \dotat{-2.89}{0} \dotat{2.74}{-4.11} \dotat{2.74}{0} \node[below] at (-2.89,0) {$-3$};
\node[above] at (2.74,0) {$3$};
\node[right] at (6.16,5.37) {$f(x)$};
}
\twochoices{$a>0;\ b=0;\ c<0;\ d>0$.}{$a>0;\ b<0;\ c<0;\ d=0$.}
{$a>0;\ b<0;\ c=0;\ d=0$.}{$a>0;\ b=0;\ c<0;\ d=0$.}
\question[TDM31]{3} Cho hàm số bậc ba $y=f(x)=ax^3+bx^2+cx+d$ có đồ thị $(C)$ như hình vẽ. Nhận xét \textbf{đúng} về dấu các hệ số $a,b,c,d$ là:
\figbox{\axes{-6.1}{6.97}{-5.56}{2.51}
\draw[grblue,thick] (-3.53,-5.04) .. controls (-2.82,-2.32) and (-2.11,0.41) .. (-1.15,0.53) .. controls (-0.19,0.66) and (1.27,-4.45) .. (2.23,-4.27) .. controls (3.19,-4.09) and (3.89,-1.21) .. (4.59,1.66);
\draw[black,thin,densely dashed] (-1.15,0)--(-1.15,0.53);
\draw[black,thin,densely dashed] (2.23,0)--(2.23,-4.27);
\dotat{-1.15}{0.53} \dotat{-1.15}{0} \dotat{2.23}{-4.27} \dotat{2.23}{0} \node[left] at (-2.63,-2.89) {$(C)$};
}
\fourchoices{$ab>0;\ cd>0$}{$ac<0;\ bd>0$}{$ad<0;\ bc<0$}{$ac>0;\ bd>0$}
\question[TDM31]{4} Cho hàm số bậc ba $y=f(x)=ax^3+bx^2+cx+d$ có đồ thị $(C)$ như hình vẽ. Nhận xét \textbf{đúng} về dấu các hệ số $a,b,c,d$ là:
\figbox{\axes{-7.27}{5.47}{-4.5}{2.97}
\draw[grblue,thick] (-5.11,-3.79) .. controls (-4.47,-1.18) and (-3.83,1.43) .. (-2.85,1.57) .. controls (-1.89,1.71) and (-0.23,-2.95) .. (0.68,-2.94) .. controls (1.59,-2.91) and (2.07,-0.61) .. (2.55,1.7);
\draw[noteborder,thin,densely dashed] (-2.85,0)--(-2.85,1.57);
\draw[noteborder,thin,densely dashed] (0.68,0)--(0.68,-2.94);
\dotat{-2.85}{1.57} \dotat{-2.85}{0} \dotat{0.68}{-2.94} \dotat{0.68}{0} \node[above] at (3.2,1.7) {$(C)$};
}
\fourchoices{$ab>0;\ cd>0$}{$ad<0;\ bc>0$}{$ac<0;\ bd>0$}{$ac>0;\ bd>0$}
\question[TDM31]{5} Cho hàm số bậc ba $y=f(x)=ax^3+bx^2+cx+d$ có đồ thị $(C)$ như hình vẽ. Nhận xét \textbf{đúng} về dấu các hệ số $a,b,c,d$ là:
\figbox{\axes{-6.3}{6.43}{-3.7}{4.23}
\draw[grblue,thick] (-3.57,2.99) .. controls (-3.01,0.41) and (-2.45,-2.17) .. (-1.76,-2.35) .. controls (-1.06,-2.53) and (-0.08,2.03) .. (0.59,1.95) .. controls (1.27,1.87) and (1.77,-0.49) .. (2.27,-2.85);
\draw[noteborder,thin,densely dashed] (-1.76,0)--(-1.76,-2.35);
\draw[noteborder,thin,densely dashed] (0.59,0)--(0.59,1.95);
\dotat{-1.76}{-2.35} \dotat{-1.76}{0} \dotat{0.59}{1.95} \dotat{0.59}{0} \node[right] at (2.43,-2.5) {$(C)$};
}
\fourchoices{$ab<0;\ cd>0$}{$ad<0;\ bc>0$}{$ac<0;\ bd<0$}{$ac>0;\ bd<0$}
\question[TDM31]{6} Cho hàm số bậc ba $y=f(x)=ax^3+bx^2+cx+d$ có đồ thị như hình vẽ. Nhận xét \textbf{đúng} về các hệ số $a,b,c,d$ là:
\figbox{\axes{-6.25}{9.02}{-4.37}{6.37}
\draw[grblue,thick] (-3.3,-3.34) .. controls (-2.99,0.37) and (-2.67,4.07) .. (-1.45,4.17) .. controls (-0.23,4.25) and (2.79,-3.05) .. (4.05,-2.77) .. controls (5.31,-2.5) and (5.71,1.65) .. (6.12,5.81);
\draw[black,thin,densely dashed] (-1.45,0)--(-1.45,4.17);
\draw[black,thin,densely dashed] (4.05,0)--(4.05,-2.77);
\dotat{-1.45}{4.17} \dotat{-1.45}{0} \dotat{4.05}{-2.77} \dotat{4.05}{0} \dotat{0}{2.8} \node[right] at (0.09,2.77) {$d$};
\node[right] at (6.12,5.81) {$f(x)$};
}
\twochoices{$a>0;\ b>0;\ c<0;\ d>0$.}{$a>0;\ b<0;\ c<0;\ d>0$.}
{$a<0;\ b<0;\ c<0;\ d>0$.}{$a>0;\ b>0;\ c>0;\ d>0$.}
\question[TDM21]{7} Đồ thị hàm số $y=x^3-3x^2+6x+1$ có tọa độ tâm đối xứng là
\fourchoices{$(0;1)$}{$(-1;-9)$}{$(1;8)$}{$(1;5)$}
\question[TDM31]{8} Gọi $(C)$ là đồ thị của hàm số $y=x^3+bx^2+cx+d$. Biết $(C)$ có tâm đối xứng là $I(2;3)$ và $y(1)=m$. Hãy tính theo $m$ giá trị của hàm số tại $x=3$.
\fourchoices{$m+1$}{$2m-3$}{$6-m$}{$m+3$}
\question[TDM31]{9} Đồ thị hàm số $y=ax^3+bx^2+cx+d$, với $a\ne0$, có hai điểm cực trị khi và chỉ khi
\fourchoices{$b^2-3ac>0$}{$b^2-4ac>0$}{$b^2-ac>0$}{$b^2-2ac>0$}
\question[TDM31]{10} Cho hàm số $y=ax^3+bx^2+cx+d$ có đồ thị như hình vẽ bên dưới. Giá trị của biểu thức $(a+2b+4c)$ bằng:
\figbox{\axes{-7.15}{6.85}{-2.73}{5.47}
\draw[grblue,thick] (-4.45,4.44) .. controls (-3.71,1.16) and (-2.97,-2.12) .. (-2.05,-2.22) .. controls (-1.15,-2.31) and (0.02,3.81) .. (0.98,3.85) .. controls (1.94,3.89) and (2.82,0.97) .. (3.7,-1.95);
\draw[noteborder,thin,densely dashed] (-2.05,0)--(-2.05,-2.22);
\draw[noteborder,thin,densely dashed] (0.98,0)--(0.98,3.85);
\dotat{-2.05}{-2.22} \dotat{-2.05}{0} \dotat{0.98}{3.85} \dotat{0.98}{0} \dotat{0}{2.37} \dotat{0}{-2.22} \node[above] at (-2.05,0) {$-2$};
\node[below] at (0.98,0) {$1$};
\node[left] at (0,2.37) {$3$};
\node[left] at (0,-2.22) {$-3$};
\node[right] at (4.12,2.94) {$f(x)$};
}
\fourchoices{$-11$}{$12$}{$16$}{$9$}
\question[TDM32]{11} Cho hàm số bậc ba $y=f(x)=ax^3+bx^2+cx+d$ có đồ thị $(C)$ và đạo hàm cấp một $f'(x)$ có đồ thị $(P)$ như hình vẽ. Khi đó giá trị của $(a+b+c+d)$ bằng
\figbox{\axes{-8.33}{5.6}{-3.47}{6.13}
\draw[grblue,thick] (-6.68,-2.77) .. controls (-5.8,0.79) and (-4.92,4.36) .. (-3.99,4.6) .. controls (-3.05,4.84) and (-1.75,-1.49) .. (-1.05,-1.33) .. controls (-0.35,-1.16) and (-0.07,2.23) .. (0.21,5.61);
\draw[red!80!black,thick,densely dotted] plot[smooth] coordinates {(-5.27,6.21) (-5.07,4.97) (-4.87,3.84) (-4.67,2.79) (-4.47,1.84) (-4.27,0.99) (-4.07,0.23) (-3.87,-0.43) (-3.67,-1) (-3.47,-1.47) (-3.27,-1.85) (-3.07,-2.14) (-2.87,-2.33) (-2.67,-2.42) (-2.47,-2.42) (-2.27,-2.33) (-2.07,-2.14) (-1.87,-1.85) (-1.67,-1.47) (-1.47,-1) (-1.27,-0.43) (-1.07,0.23) (-0.87,0.99) (-0.67,1.84) (-0.47,2.79) (-0.27,3.84) (-0.07,4.97) (0.13,6.21)};
\draw[black,thin,densely dashed] (-3.99,4.6)--(-3.99,0);
\draw[black,thin,densely dashed] (-2.57,0)--(-2.57,-2.43);
\draw[black,thin,densely dashed] (-2.57,-2.43)--(-0.07,-2.43);
\draw[black,thin,densely dashed] (-1.05,0)--(-1.05,-1.33);
\dotat{-3.99}{4.6} \dotat{-3.99}{0} \dotat{-2.57}{0} \dotat{-2.57}{-2.43} \dotat{-1.05}{-1.33} \dotat{-1.05}{0} \dotat{0}{0} \dotat{0}{3.43} \node[below left] at (-4,0) {$-3$};
\node[above] at (-2.57,0) {$-2$};
\node[right] at (-0.07,-2.43) {$-3$};
\node[left] at (-6.4,4.4) {$(P)$};
\node[left] at (-6.47,-2.67) {$(C)$};
}
\fourchoices{$20$}{$0$}{$18$}{$25$}
\question[TDM32]{12} Cho hàm số bậc ba $y=f(x)$ và đạo hàm cấp hai $f''(x)$ có đồ thị lần lượt là $(C)$ và $\Delta$ như hình vẽ. Nhận xét \textbf{đúng} về giá trị $f(3)$ là
\figbox{\axes{-8.31}{8.56}{-4.32}{4.61}
\draw[grblue,thick] (-6.47,-2.68) .. controls (-5.44,0.4) and (-4.41,3.49) .. (-3.11,3.4) .. controls (-1.83,3.31) and (0,-3.16) .. (1.27,-3.23) .. controls (2.53,-3.29) and (3.51,-0.13) .. (4.49,3.03);
\draw[black,thin] (-7.74,-2.31)--(6.93,2.61);
\draw[black,thin,densely dashed] (0,1.69)--(4.09,1.69);
\draw[black,thin,densely dashed] (4.09,1.69)--(4.09,0);
\dotat{-0.87}{0} \dotat{0}{1.69} \dotat{4.09}{1.69} \dotat{4.09}{0} \node[below] at (-0.87,0) {$-1$};
\node[below] at (4.09,0) {$4$};
\node[left] at (0,1.69) {$30$};
\node[above] at (6.96,2.61) {$\Delta$};
\node[above] at (4.49,3.08) {$(C)$};
}
\twochoices{$-15<f(3)<-10$}{$25<f(3)<30$}{$15<f(3)<25$}{$-10<f(3)<0$}
\parttwo{13}{16}
\question[TDM31]{13} Cho hàm số bậc ba $y=f(x)=ax^3+bx^2+cx+d$ có đồ thị $(C)$ như hình vẽ. Hỏi mệnh đề nào dưới đây là đúng hay sai?
\figbox{\axes{-4.59}{5.68}{-4.28}{3.72}
\draw[grblue,thick] (-3.24,2.67) .. controls (-2.35,0.06) and (-1.46,-2.55) .. (-0.61,-2.78) .. controls (0.24,-3.02) and (1.08,1.42) .. (1.89,1.27) .. controls (2.7,1.11) and (3.47,-1.31) .. (4.24,-3.72);
\dotat{0}{0} \node[right] at (2.48,1.72) {$(C)$};
}
\tfbox{Số thực $a<0$.}
{Đồ thị cắt trục tung tại điểm có tung độ dương.}
{Phương trình $f(x)=0$ có đúng một nghiệm thực dương.}
{Giá trị $bc>0$.}
\question[TDM31]{14} Cho hàm số bậc ba $y=f(x)=x^3+bx^2+cx+d$ và đạo hàm cấp hai $f''(x)$ có đồ thị lần lượt là $(C)$ và $\Delta$ như hình vẽ. Hỏi mệnh đề nào dưới đây là đúng hay sai?
\figbox{\axes{-7.02}{7.11}{-4.56}{3.31}
\draw[grblue,thick] (-6.73,-3.84) .. controls (-5.7,-0.75) and (-4.66,2.33) .. (-3.37,2.25) .. controls (-2.09,2.16) and (-0.25,-4.32) .. (1.01,-4.39) .. controls (2.27,-4.45) and (3.25,-1.29) .. (4.23,1.86);
\draw[black,thin] (-7.95,-3.19)--(5.91,3.25);
\draw[black,thin,densely dashed] (-3.37,2.25)--(0,2.25);
\draw[black,thin,densely dashed] (1.01,0)--(1.01,-4.39);
\dotat{-3.37}{2.25} \dotat{0}{2.25} \dotat{0}{0} \dotat{-1.16}{0} \dotat{1.01}{-4.39} \dotat{1.01}{0} \node[above left] at (0,2.25) {$2$};
\node[above] at (-1.16,0) {$-1$};
\node[above] at (1.01,0) {$1$};
\node[above] at (5.78,3.24) {$\Delta$};
\node[right] at (3.38,-2.89) {$(C)$};
}
\tfbox{Số thực $b=3$.}
{Số thực $c=9$.}
{Số thực $d=-20$.}
{Giá trị biểu thức $b+c+d=-26$.}
\question[TDM31]{15} Cho hàm số bậc ba $y=f(x)=ax^3+bx^2+cx+d$ có đồ thị $(C)$ như hình vẽ. Nhận xét \textbf{đúng} về giá trị thực của tham số $m$ là gì? Hỏi mệnh đề nào dưới đây là đúng hay sai?
\figbox{\axes{-9.4}{7.9}{-5.75}{6.3}
\draw[grblue,thick] (-5.6,5.65) .. controls (-4.8,1.9) and (-3.95,-3.1) .. (-2.85,-3.1) .. controls (-1.75,-3.1) and (-1.05,3.05) .. (0,3.05) .. controls (1.25,3.05) and (2.65,-2.35) .. (3.5,-4.4);
\draw[black,thin,densely dashed] (-2.85,0)--(-2.85,-3.1);
\draw[black,thin,densely dashed] (-2.85,-3.1)--(0,-3.1);
\node[above] at (-2.85,0) {$m$};
\node[above] at (2.2,0) {$1$};
\node[above] at (0,3.05) {$2$};
\node[right] at (0,-3.1) {$-2$};
\node[right] at (3.7,-4.4) {$(C)$};
}
\tfbox{$d=2$.}{$c=0$.}{Giá trị nhỏ nhất của hàm số $f(x)$ trên đoạn $[m;1]$ bằng $-2$.}{Giá trị $m=-2$.}
\question[TDM31]{16} Cho hàm số bậc ba $y=f(x)=ax^3+bx^2+cx+d$ có đồ thị $(C)$. Biết $I(1;2)$ là tâm đối xứng của $(C)$ và $A(-1;18)$ là một điểm cực trị của $(C)$. Hỏi mệnh đề nào dưới đây là đúng hay sai?
\tfbox{Điểm cực trị khác của $(C)$ có tọa độ là $(3;-11)$.}
{$a>0$.}{$b<0$.}{$a+2b+c+3d=25$.}
\partthree{17}{22}
\question[TDM31]{17} Biết đồ thị hàm số $y=x^3-6x^2+2x+10$ có tâm đối xứng là $I(a;b)$. Hãy tính giá trị biểu thức $(a+b)$?
\shortanswer
\question[TDM31]{18} Biết đồ thị của hàm số $y=f(x)=x^3+bx^2+cx+d$ có một điểm cực trị là $A(1;2)$. Khi đó giá trị biểu thức $T=b+2c+3d$ bằng bao nhiêu?
\shortanswer
\question[TDM31]{19} Cho hàm số bậc ba có đồ thị $(C)$ được biểu diễn như hình vẽ bên dưới. Khi đó giá trị của tổng $(a+b)$ bằng bao nhiêu?
\figbox{\axes{-9.2}{8.4}{-5.6}{5.15}
\draw[grblue,thick] (-7.05,-4) .. controls (-5.8,0.4) and (-4.5,3.65) .. (-2.95,3.65) .. controls (-1.7,3.65) and (-0.1,-3.3) .. (1.1,-3.3) .. controls (2.6,-3.3) and (3.8,0.95) .. (4.7,3.75);
\draw[black,thin,densely dashed] (-2.95,3.65)--(-2.95,0);
\draw[black,thin,densely dashed] (1.1,0)--(1.1,-3.3);
\dotat{-2.95}{3.65} \dotat{-2.95}{0} \dotat{1.1}{0} \dotat{1.1}{-3.3} \node[below] at (-5.45,0) {$-6$};
\node[below] at (-2.95,0) {$a$};
\node[above] at (-1.2,0) {$-1$};
\node[above] at (1.1,0) {$b$};
\node[below] at (3.75,0) {$4$};
\node[right] at (5.1,3.4) {$(C)$};
}
\shortanswer
\question[TDM32]{20} Cho hàm số $y=(m-1)x^3-(m-3)x-2m+2022$. Gọi $S$ là tập tất cả các giá trị nguyên của tham số $m$ để hàm số đồng biến trên $\mathbb R$. Số phần tử tập $S$ bằng bao nhiêu?
\shortanswer
\question[TDM32]{21} Cho hàm số bậc ba $y=x^3-3mx^2+2m$ có đồ thị $(C)$. Gọi điểm cực đại và điểm cực tiểu lần lượt là $A$ và $B$. Để tam giác $OAB$ cân tại $O$ thì giá trị $m^2$ bằng bao nhiêu? (Kết quả lấy một chữ số sau dấu phẩy.)
\shortanswer
\question[TDM42]{22} Cho hàm số bậc ba $y=x^3-3(m+1)x^2+(6m^2+12m)x$. Hàm số đạt cực trị tại $x_1$ và $x_2$. Hãy tìm giá trị nhỏ nhất của biểu thức $T=(x_1^2-5)(1-x_2^2)$?
\shortanswer
\finishExam
\clearpage
\examtitle{TDMXX11}{KHẢO SÁT VÀ VẼ ĐỒ THỊ HÀM PHÂN THỨC}
\partone{10}
\question[TDM21]{1} Đồ thị hàm số $y=\dfrac{2x-3}{x+1}$ có tâm đối xứng là
\fourchoices{$A(1;2)$}{$B(-1;2)$}{$C(-1;-2)$}{$D(-1;3)$}
\question[TDM31]{2} Đồ thị hàm số $y=\dfrac{2x-3}{x+1}$ có trục đối xứng (không cắt đồ thị) có phương trình là
\fourchoices{$y=x+3$}{$y=-x+1$}{$y=2x-1$}{$y=x+1$}
\question[TDM31]{3} Cho hàm số $f(x)=\dfrac{ax+b}{cx+d}$ có đồ thị cho như hình vẽ. Nhận xét đúng khi nói về dấu các hệ số $a,b,c,d$ là:
\figbox{\axes{-4.03}{6.8}{-5.63}{3.97}
\draw[black,thick] (-4.34,-0.71) .. controls (-2.83,-0.69) and (-2,-0.43) .. (-1.49,0) .. controls (-0.74,0.37) and (0.49,1.91) .. (0.63,3.29);
\draw[black,thick] (1.49,-5.4) .. controls (1.6,-3.4) and (3.6,-1.51) .. (6.54,-1.4);
\draw[black,thin] (1.06,3.31)--(1.06,-5.63);
\draw[black,thin] (-4.34,-1.06)--(6.54,-1.06);
}
\twochoices{$\begin{cases}ab>0\\cd>0\end{cases}$}{$\begin{cases}ad>0\\bc<0\end{cases}$}
{$\begin{cases}ad<0\\bc>0\end{cases}$}{$\begin{cases}ac>0\\bd>0\end{cases}$}
\question[TDM31]{4} Cho hàm số $f(x)=\dfrac{ax+b}{cx+d}$ có đồ thị cho như hình vẽ. Nhận xét đúng khi nói về dấu các hệ số $a,b,c,d$ là:
\figbox{\axes{-2.96}{4.98}{-4.73}{2.88}
\draw[black,thick] (-2.69,-0.96) .. controls (-1.44,-0.96) and (-0.02,-2) .. (0.29,-3.12) .. controls (0.5,-3.71) and (0.54,-4.25) .. (0.54,-4.5);
\draw[black,thick] (0.92,2.58) .. controls (0.94,1.29) and (1.9,0.08) .. (2.56,0) .. controls (3.48,-0.54) and (4.31,-0.65) .. (4.92,-0.6);
\draw[black,thin] (0.77,2.83)--(0.77,-4.73);
\draw[black,thin] (-2.96,-0.77)--(4.98,-0.77);
}
\twochoices{$\begin{cases}ab<0\\cd>0\end{cases}$}{$\begin{cases}ab>0\\cd<0\end{cases}$}
{$\begin{cases}ad>0\\bc>0\end{cases}$}{$\begin{cases}ac>0\\bd<0\end{cases}$}
\question[TDM31]{5} Cho hàm số $y=\dfrac{ax+b}{cx+d}$ có đồ thị $(C)$ và nhận điểm $(2;-1)$ làm tâm đối xứng. Đồ thị $(C)$ cắt hai trục tọa độ lần lượt tại $M$ và $N$ sao cho diện tích tam giác $OMN$ bằng $2$. Hàm số đã cho tương ứng là:
\twochoices{$y=\dfrac{-x+2\sqrt2}{x-2}$}{$y=\dfrac{x-1}{x-2}$}
{$y=\dfrac{x+2\sqrt2}{x-2}$}{$y=\dfrac{x+2}{-x+2}$}
\question[TDM31]{6} Cho hàm số $y=f(x)=\dfrac{3x-2}{x-1}$ có đồ thị $(C)$. Hãy tìm tất cả các giá trị thực của tham số $m$ để phương trình $f(x)=2m-1$ vô nghiệm?
\fourchoices{$(-\infty;1)$}{$m=2$}{$m=\dfrac12$}{$m\in(3;+\infty)$}
\question[TDM31]{7} Đồ thị hàm số $y=\dfrac{x^2-5x+7}{x-2}$ có tâm đối xứng là
\fourchoices{$(2;3)$}{$(-2;-5)$}{$(2;-1)$}{$(1;-2)$}
\question[TDM31]{8} Biết đồ thị hàm số $y=\dfrac{x^2-3x+4}{x-1}$ có trục đối xứng là đường thẳng có phương trình $y=ax+b$, với $a>0$. Giá trị biểu thức $T=a+b$ bằng
\fourchoices{$2$}{$3$}{$0$}{$-1$}
\question[TDM31]{9} Hỏi đồ thị hàm số $y=\dfrac{2x+4}{x-1}$ có tất cả bao nhiêu điểm có tọa độ nguyên?
\fourchoices{$8$}{$9$}{$7$}{$6$}
\question[TDM31]{10} Hỏi đồ thị hàm số $y=\dfrac{x^2-5x+10}{x-2}$ có tất cả bao nhiêu điểm có tọa độ nguyên?
\fourchoices{$8$}{$4$}{$2$}{$6$}
\parttwo{11}{12}
\question[TDM31]{11} Cho hàm số $y=f(x)=\dfrac{x^2-4x+3}{x-2}$ có đồ thị $(C)$. Hỏi mệnh đề nào dưới đây là đúng hay sai?
\tfbox{Đường tiệm cận đứng của $(C)$ là $x=2$ và đường tiệm cận xiên của $(C)$ là $y=x-2$.}
{Hàm số $f(x)$ có hai điểm cực trị.}
{Đồ thị $(C)$ có tâm đối xứng là $I(2;0)$.}
{Một trong hai trục đối xứng của $(C)$ có phương trình là $y=(\sqrt2+1)x-2\sqrt2-2$.}
\question[TDM32]{12} Cho hàm số $y=\dfrac{2x-11}{x-1}$ có đồ thị $(C)$. Hỏi mệnh đề nào dưới đây là đúng hay sai?
\tfbox{Tâm đối xứng của $(C)$ là $I(1;2)$.}
{Đồ thị $(C)$ giao với $Ox$ tại $A$ và giao với $Oy$ tại $B$. Có $AB=6$.}
{Số điểm có tọa độ nguyên thuộc đồ thị $(C)$ là $8$.}
{Số parabol giao với $(C)$ tại ba điểm phân biệt có tọa độ nguyên là $20$.}
\partthree{13}{17}
\question[TDM31]{13} Cho đồ thị $(C)$: $y=\dfrac{x-3}{x+1}$. Biết rằng có hai điểm phân biệt thuộc đồ thị $(C)$ và cách đều hai trục tọa độ. Giả sử các điểm đó lần lượt là $M$ và $N$. Hãy xác định độ dài đoạn thẳng $MN$ (làm tròn kết quả đến hàng phần trăm)?
\shortanswer
\question[TDM32]{14} Cho hàm số $y=\dfrac{x+3}{x-1}$ có đồ thị $(C)$ và đường thẳng $d:y=x-2m$. Để $(C)$ cắt $d$ tại hai điểm $A,B$ phân biệt sao cho $AB=2\sqrt{10}$ thì giá trị $m^2$ bằng bao nhiêu?
\shortanswer
\question[TDM32]{15} Cho hàm số $y=\dfrac{2x-1}{x-2}$ có đồ thị $(C)$ và đường thẳng $d:y=x+m$. Để $(C)$ cắt $d$ tại hai điểm phân biệt $A,B$, hỏi độ dài $AB$ ngắn nhất bằng bao nhiêu (làm tròn kết quả đến hàng phần mười)?
\shortanswer
\question[TDM32]{16} Cho hàm số $y=\dfrac{x-1}{x+1}$ có đồ thị $(C)$. Biết $(C)$ cắt trục $Ox$ và $Oy$ lần lượt tại hai điểm $A$ và $B$. Gọi $M(a;b)$ là một điểm nằm trên $(C)$ sao cho diện tích tam giác $MAB$ bằng $3$. Hỏi có tất cả bao nhiêu điểm $M$ thỏa mãn bài toán?
\shortanswer
\question[TDM32]{17} Cho hàm số $y=f(x)=\dfrac{3x-2}{x+1}$ có đồ thị $(C)$. Hỏi có tất cả bao nhiêu đường thẳng cắt $(C)$ tại hai điểm phân biệt có tọa độ nguyên?
\shortanswer
\finishExam
% Nguon: ba de TDMXX09, TDMXX10, TDMXX11 (Nguyen Dang Ai -- Tu Duy Mo).
% Ban go LaTeX: BTVN_B5.tex theo phong cach BTVN_01_B1.tex; do thi ve lai TikZ.
\end{document}